\documentclass[10pt,a4paper]{amsart}
\usepackage[margin=19mm]{geometry}
\usepackage{amsmath,amssymb,mathtools}
\usepackage[scr=rsfs,cal=euler]{mathalfa}
\usepackage{xfrac}
\usepackage[unicode,hidelinks,bookmarks=false]{hyperref}
\newcommand{\ie}{i.e.\ }
\newcommand{\cf}{cf.\ }
\newcommand{\resp}{respectively\ }
\newcommand{\Subsection}{\subsection}
\providecommand{\nobreakdash}{\nobreak\hbox{-}}
\setlength{\emergencystretch}{2em}
\multlinegap0pt
\usepackage[shortlabels]{enumitem}
\usepackage{mathtools}
\usepackage{amssymb}
\usepackage{xcolor}
\usepackage{algorithm}\usepackage{algpseudocode}
\newtheorem{definition}{Definition}
\newtheorem{assumption}{Assumption}
\newtheorem{theorem}{Theorem}
\newcommand{\SOL}{\mathrm{SOL}}
\title{Complete Trajectory Tracking for Polynomial Differential Variational Inequalities: A Moment-SOS Based Structure-Aware Approach}
\author{Zhang Huang, Wei Li, Jie Wang}
\date{Source: arXiv:2606.09609v1 (CC BY 4.0)}
\begin{document}
\maketitle

\noindent\emph{Source and licence.} Complete Trajectory Tracking for Polynomial Differential Variational Inequalities: A Moment-SOS Based Structure-Aware Approach. Version arXiv:2606.09609v1 (CC BY 4.0), distributed by arXiv under the Creative Commons Attribution 4.0 International licence, http://creativecommons.org/licenses/by/4.0/, as recorded in the arXiv licence metadata for this exact version and shown on the abstract page https://arxiv.org/abs/2606.09609v1. Full licence notice: \texttt{licenses/arxiv-2606.09609-CC-BY-4.0.txt}.\par\smallskip

\noindent\textit{Editorial scope.} Counted unit: the article's theorem thm:maximal\_existence (source0.tex lines 563--578). The article's complete original proof of it is delivered with it (source0.tex lines 580--608, one \texttt{proof} environment of 29 lines). No mathematics is written, repaired or re-derived: the statement and the proof are the article's own lines, in the article's own notation, and no retained line is substituted or re-flowed.  Retained uncounted as the source-local context the proof needs: lines 117--134; lines 220--229; lines 243--245; lines 308--314.  Nothing was omitted from any retained span, so the package carries no omission record.  No retained line contains a citation, so the wrapper carries no bibliography and no bibliography entry.  No retained line contains a figure environment, an image-inclusion command or drawing code, so no image is delivered and the package declares no asset dependency.  Editorial formatting only, in addition: an AMS \texttt{amsart} preamble that restates the article's own declarations and macros used by the retained lines, namely the environments \texttt{definition}, \texttt{assumption}, \texttt{theorem} on separate counters, together with the standard packages those lines need.  The retained environments print in this document as Definition 1, Assumption 1, Theorem 1 in order of appearance, whereas the article prints the same items with the numbering of its own full document.  The complete native source of the article as distributed in the arXiv e-print for this exact version is bundled unchanged as \texttt{sources/202412/source0.tex}.
\begin{definition}
\label{def:pdvi}
Let $T > 0$ be a fixed time horizon and $x_0 \in \mathbb{R}^n$ an initial state. Consider the system over $t \in [0, T]$ given by
\begin{subequations}\label{eq:pdvi-system}
\begin{align}
    \dot{x}(t) &= f\big(x(t), u(t)\big), \label{eq:pdvi-dyn} \\
    \langle F(x(t), u(t)),\, z - u(t) \rangle &\geq 0 \quad \text{for all } z \in K, \label{eq:pdvi-vi} \\
    x(0) &= x_0. \label{eq:pdvi-init}
\end{align}
\end{subequations}
This system is called a \emph{Polynomial Differential Variational Inequality} (PDVI) if the following structural conditions are satisfied: the vector field $f : \mathbb{R}^n \times \mathbb{R}^m \to \mathbb{R}^n$ and the variational inequality mapping $F : \mathbb{R}^n \times \mathbb{R}^m \to \mathbb{R}^m$ are polynomial mappings, i.e., $f \in \mathbb{R}[x, u]^n$ and $F \in \mathbb{R}[x, u]^m$, and the constraint set $K \subset \mathbb{R}^m$ is a basic closed semi-algebraic set of the form
\[
    K = \left\{ u \in \mathbb{R}^m \;\middle|\; g_i(u) \geq 0,\ i = 1, \dots, p \right\},
\]
where $p \in \mathbb{N}$ and each $g_i \in \mathbb{R}[u]$ is a real polynomial.

A pair of functions $(x(\cdot), u(\cdot))$ is said to be a \emph{Carath\'{e}odory solution} of the PDVI if $x : [0, T] \to \mathbb{R}^n$ is absolutely continuous, $u : [0, T] \to \mathbb{R}^m$ is Lebesgue measurable and essentially bounded, and the relations~\eqref{eq:pdvi-dyn}--\eqref{eq:pdvi-init} hold for almost every $t \in [0, T]$.
\end{definition}

\begin{assumption} \label{ass:core}
We assume that the following conditions hold for the parametric differential variational inequality (PDVI) specified in Definition~\ref{def:pdvi}.
\begin{enumerate}[label=\textup{(A\arabic*)}, widest=(A3), left=0pt, align=left]
    \item The control constraint set $K$ is nonempty, bounded, closed, and basic semi-algebraic.

    \item The linear independence constraint qualification (LICQ) holds at every point in $K$.

    \item The mappings $F \colon \mathbb{R}^n \times K \to \mathbb{R}^m$ and $f \colon \mathbb{R}^n \times K \to \mathbb{R}^n$ are polynomial in their arguments.
\end{enumerate}
\end{assumption}

\begin{theorem} \label{thm:eq-semi-graph}
Under Assumption~\ref{ass:core}, the graph $\mathrm{Graph}(\mathcal{S})$ is a semi-algebraic subset of $\mathbb{R}^{n+m}$.
\end{theorem}

\begin{theorem} \label{thm:selection}
Under Assumption~\ref{ass:core}, let $\mathcal{S}: \mathcal{D} \rightrightarrows K$ be the solution mapping of the parameterized polynomial variational inequality, where $\mathcal{D} := \{x \in \mathbb{R}^n \mid \mathcal{S}(x) \neq \emptyset\}$ and $K \subset \mathbb{R}^m$ is compact. Then
\begin{enumerate}
    \item There exists a Lebesgue null set $\mathcal{N} \subset \mathcal{D}$ such that for any $x_0 \in \mathcal{D} \setminus \mathcal{N}$ and any $u_0 \in \mathcal{S}(x_0)$, there exists an open neighborhood $U \subset \mathcal{D} \setminus \mathcal{N}$ of $x_0$ and a continuous (indeed, semi-algebraic) function $\sigma: U \to K$ satisfying $\sigma(x_0) = u_0$ and $\sigma(x) \in \mathcal{S}(x)$ for all $x \in U$.
    \item There exists a Borel measurable function $\sigma: \mathcal{D} \to K$ such that $\sigma(x) \in \mathcal{S}(x)$ for all $x \in \mathcal{D}$.
\end{enumerate}
\end{theorem}

\begin{theorem} \label{thm:maximal_existence}
Consider the polynomial differential variational inequality (PDVI) system
\begin{subequations}
\begin{align}
\dot{x}(t) &= f(x(t), u(t)), \quad \text{a.e. } t \in [0, T), \label{eq:pdvi_dyn_max} \\
u(t) &\in \SOL(K, F(x(t), \cdot)), \quad \forall t \in [0, T), \label{eq:pdvi_vi_max} \\
x(0) &= x_0, \label{eq:pdvi_ic_max}
\end{align}
\end{subequations}
under Assumption~\ref{ass:core}. Let $\mathcal{S}(x) := \SOL(K, F(x, \cdot))$ and $\mathcal{D} := \{x \in \mathbb{R}^n : \mathcal{S}(x) \neq \emptyset\}$. Then, for any initial condition $x_0 \in \mathcal{D}$, there exists a unique maximal time $T_{\max} \in (0, \infty]$ and a Carath\'{e}odory solution $(x(\cdot), u(\cdot))$ to the PDVI system \eqref{eq:pdvi_dyn_max}--\eqref{eq:pdvi_ic_max} defined on the maximal interval $[0, T_{\max})$. Furthermore, if $T_{\max} < \infty$, then the trajectory $x(\cdot)$ leaves every compact subset of $\mathcal{D}$ as $t \nearrow T_{\max}$; that is,
\[ 
\lim_{t \nearrow T_{\max}} \mathrm{dist}\big(x(t), \partial \mathcal{D}\big) = 0 
\quad \text{or} \quad 
\lim_{t \nearrow T_{\max}} \|x(t)\| = \infty.
\]
\end{theorem}

\begin{proof}
By Theorem~\ref{thm:eq-semi-graph}, the graph $\mathrm{Graph}(\mathcal{S})$ is a semi-algebraic subset of $\mathbb{R}^{n+m}$. Since $K$ is compact and $F$ is polynomial, each fiber $\mathcal{S}(x)$ is a nonempty compact semi-algebraic set for $x \in \mathcal{D}$, and $\mathcal{D} = \pi_x(\mathrm{Graph}(\mathcal{S}))$ is semi-algebraic by the Tarski--Seidenberg theorem.

By Theorem~\ref{thm:selection}(2), there exists a Borel measurable selection $\sigma : \mathcal{D} \to K$ such that $\sigma(x) \in \mathcal{S}(x)$ for all $x \in \mathcal{D}$. Consider the ODE
\[
\dot{x}(t) = f(x(t), \sigma(x(t))), \quad x(0) = x_0.
\]
Since $f$ is a polynomial (hence continuous) and $\sigma$ is Borel measurable, the right-hand side $t \mapsto f(x(t), \sigma(x(t)))$ is measurable whenever $x(\cdot)$ is measurable, and locally bounded because $f$ is continuous and both $x(t)$ and $\sigma(x(t))$ remain in compact sets over finite time intervals (as shown below). Therefore, by Carath\'{e}odory's existence theorem, there exists $\epsilon > 0$ and an absolutely continuous function $x : [0, \epsilon) \to \mathbb{R}^n$ satisfying the ODE almost everywhere with $x(0) = x_0$. Setting $u(t) := \sigma(x(t))$, we obtain a Carath\'{e}odory solution $(x(\cdot), u(\cdot))$ to the PDVI on $[0, \epsilon)$.

Let $[0, T_{\max})$ be the maximal interval of existence of this solution. Suppose, for contradiction, that $T_{\max} < \infty$ but that the trajectory remains in some compact set $C \subset \mathcal{D}$ for all $t \in [0, T_{\max})$. Since $C$ is compact and contained in the semi-algebraic set $\mathcal{D}$, and since $f$ is continuous and $K$ is compact, the set
\[
\{ f(x, u) : x \in C,\, u \in K \}
\]
is bounded. Hence $\|\dot{x}(t)\|$ is essentially bounded on $[0, T_{\max})$, implying that $x(\cdot)$ is uniformly Lipschitz continuous on this interval. Consequently, the limit
\[
x(T_{\max}^-) := \lim_{t \nearrow T_{\max}} x(t)
\]
exists in $\mathbb{R}^n$. Because $C$ is compact and $C \subset \mathcal{D}$, we have $x(T_{\max}^-) \in C \subset \mathcal{D}$.

Now, since $x(T_{\max}^-) \in \mathcal{D}$, the solution set $\mathcal{S}(x(T_{\max}^-))$ is nonempty. By Theorem~\ref{thm:selection}(2) again, there exists a Borel measurable selection $\sigma$ defined in a neighborhood of $x(T_{\max}^-)$. Applying Carath\'{e}odory's theorem with initial condition $x(T_{\max}^-)$ yields a local solution on $[T_{\max}, T_{\max} + \delta)$ for some $\delta > 0$. Concatenating this with the original solution contradicts the maximality of $T_{\max}$.

Therefore, if $T_{\max} < \infty$, the trajectory $x(t)$ cannot remain in any compact subset of $\mathcal{D}$. This implies that either $\|x(t)\| \to \infty$ or $x(t)$ approaches the boundary of $\mathcal{D}$ as $t \nearrow T_{\max}$. Since $\mathcal{D}$ is semi-algebraic, its topological boundary coincides (up to a lower-dimensional set) with $\partial \mathcal{D}$, and thus
\[
\lim_{t \nearrow T_{\max}} \mathrm{dist}(x(t), \partial \mathcal{D}) = 0
\quad \text{or} \quad
\lim_{t \nearrow T_{\max}} \|x(t)\| = \infty,
\]
as claimed.
\end{proof}

\end{document}
